
% contents of the chapter Quantum mechanics
\chapter{Quantum Mechanics}
\label{quant-mech}
{
    \setlength\epigraphwidth{0.6\textwidth}
\epigraph{ Quantum mechanics is a
different kind of a theory to represent the world.
} {R. P. Feynman}
}
The relation between classical and quantum mechanics is often
described as one of extension: quantum mechanics generalizes
classical mechanics. Actually it is correct only for a
part (though very important one) of classical
 systems ---  systems with phase space $ \mathbb{R}^{2 n} $
 and  Hamiltonian which is quadratic function of momenta,
 and beside them some relatively exotic highly symmetric systems. 
 In that case quantum mechanics
replaces classical mechanics by altering its fundamental
mathematical structure. Classical mechanics reappears only 
as some kind of  singular limit.
On the other hand  crucial components of classical theory 
like constraints, canonical transformations, trajectories,
have no faithful quantum analogues.

Non-relativistic quantum mechanics can be presented axiomatically,
without explicit reference to a classical precursor. 
However all rigorous definitions and theorems give us no way to connect 
the mathematical constructions to the real world. 
"Of course, nature knows the quantum mechanics,
and the classical mechanics is only an approximation; so there is no mystery in
the fact that in classical mechanics there is some shadow of quantum mechanical
laws---which are truly the ones underneath. To reconstruct the original object
from the shadow is not possible in any direct way, but the shadow does help
you to remember what the object looks like" \cite{flp3}. 
So the problem of \emph{quantization}---the transition from 
classical to quantum descriptions---actually works backward.

\section{Bra-Ket Notation}
\label{bra-ket-notation}

Bra--ket notation is the standard symbolic language of quantum
mechanics. While extraordinarily powerful and conceptually elegant,
it compresses multiple layers of mathematical structure into compact
expressions. 
Its success stems from its ability to encode:
\begin{itemize}
\item Linear algebraic structure,
\item Operator theory,
\item Spectral decomposition,
\item Measurement theory,
\item Tensor products and entanglement.
\end{itemize}
However, the notation often suppresses functional-analytic and
geometric subtleties. 

It appears at Paul Dirac's book \cite{dirac-pqm}  "The Principles of Quantum Mechanics" (first edition 1930). Let's describe pros and cons:
\subsection{ Pros}
\begin{enumerate}
   \item
 Coordinate-free and basis-independent
bra--ket notation expresses quantum states abstractly:
\[|\psi\rangle \in \mathcal{H}, \qquad \langle \phi| \in \mathcal{H}^*\]
This makes the formalism:
\begin{itemize}
\item
  independent of representation (position, momentum, spin, etc.)
\item
  cleanly adaptable to changes of basis
\item
  naturally compatible with symmetry arguments
\end{itemize}
Example:
\[\psi(x) = \langle x|\psi\rangle\]
The same state appears as a function only after choosing the position
basis.
\item
 Clear dual structure
  The notation explicitly encodes:
\begin{itemize}
\item
  kets: vectors \(|\psi\rangle\)
\item
  bras: dual vectors \(\langle\psi|\)
\item
  inner product:
\end{itemize}
\[\langle \phi | \psi \rangle\]
This makes conjugation and adjoints transparent:
\[(|\psi\rangle)^\dagger = \langle \psi|\]
\item
elegant operator representation:
\[A|\psi\rangle\]
outer products generate rank-one operators:
\[|\psi\rangle\langle\phi|\]
projection operators:
\[P = |a\rangle\langle a|\]
spectral decomposition:
\[A = \sum_a a |a\rangle\langle a|\]
or (continuous spectrum)
\[A = \int a\, |a\rangle\langle a|\, da\]
This makes the spectral theorem visually intuitive.
\item
Simplifies tensor products ---
composite systems are naturally written:
\[|\psi\rangle \otimes |\phi\rangle\]
Entangled states appear transparently:
\[|\Psi\rangle = \sum_{ij} c_{ij} |i\rangle \otimes |j\rangle\]
\item
Works naturally with PVM and POVM formalism
projection-valued measure:
\[E(\Delta) = \int_\Delta |x\rangle\langle x| dx\]
POVM elements:
\[F_i = M_i^\dagger M_i\]
Bra--ket notation makes these constructions compact and readable.
\end{enumerate}

\subsection{ Cons}
\begin{enumerate}
   \item
Lack of mathematical rigor (without extra care)---
Dirac freely wrote expressions like:
\[\langle x|x' \rangle = \delta(x-x').\]
But:
\begin{itemize}
\item
  \(|x\rangle \notin L^2(\mathbb{R})\)
\item
  delta functions are distributions, not functions
\item
  continuous ``eigenvectors'' are not true Hilbert space elements.
\end{itemize}
Strictly speaking, this requires the theory of rigged Hilbert spaces:
\[\Phi \subset \mathcal{H} \subset \Phi',\]
Without this, the notation hides subtle functional analysis.
   \item
Can hide domain issues ---
unbounded operators (like momentum or Hamiltonians) require domain care:
\[P = -i\hbar \frac{d}{dx}.\]
In bra--ket notation, one may forget:
\begin{itemize}
\item
  domain restrictions
\item
  self-adjoint extensions
\item
  boundary conditions.
\end{itemize}
These are essential in rigorous analysis.
   \item
Ambiguity in infinite dimensions---
in finite dimensions:
\[\langle \psi| = |\psi\rangle^\dagger\]
in infinite dimensions:
\begin{itemize}
\item
  The identification depends on Riesz representation.
\item
  Not all linear functionals are representable without topology
  assumptions.
\end{itemize}
The notation may conceal these subtleties.
   \item
 Overuse leads to symbolic manipulation without understanding.
Some common problems:
\begin{itemize}
\item
  Treating kets like column vectors everywhere.
\item
  Confusing formal identities with rigorous equalities.
\item
  Ignoring operator domains and closure.
\end{itemize}
It encourages formal calculation over structural analysis.
\end{enumerate}

\subsection{Some Explanations}
\label{some-explanations}

Dirac's book \cite{dirac-pqm} contains
statements that are mathematically unclear, incomplete, or (by modern standards) incorrect. 
But he did not work in the modern functional-analytic framework that developed later.
For example:
\begin{itemize}
   \item He freely used “eigenvectors” of operators with continuous spectrum.
    \item He manipulated $ \delta $-functions before distribution theory existed.
    \item He treated unbounded operators without domain analysis.
    \item He wrote symbolic identities like \[ \langle x|p \rangle=\exp(i p x/\hbar) \]
as if everything were finite-dimensional.
\end{itemize}
From today’s perspective, this is not rigorous.
From 1930’s perspective, it was revolutionary.

When Dirac wrote his book:
\begin{itemize}
   \item Hilbert space theory was still developing.
\item Spectral theorem was not yet standard physics language.
\item Rigged Hilbert spaces did not exist.
\item Distribution theory did not exist.
\end{itemize}
He was not misusing established mathematics.
He was inventing a new symbolic calculus before the mathematics to justify it existed.
Later, mathematicians 
provided rigorous frameworks that made Dirac’s manipulations meaningful.

However some of statements are genuinely incorrect, for example:
\begin{itemize}
   \item implicit assumption that every observable has a complete set of eigenvectors;
\item overuse of spectral decompositions without domain care;
\item loose treatment of continuous spectra;
\item ambiguities about self-adjoint vs symmetric operators;
\item (most serious!) statement that canonical transformations in classical mechanics 
    may be lifted to unitary transformations in quantum mechanics.
\end{itemize}

In his book \cite{neumann-mfqm}  Mathematische Grundlagen der Quantenmechanik (1932) von Neumann made everything measure-theoretic and Hilbert-space rigorous and avoided $ \delta $-functions entirely.
Though von Neumann’s formulation is mathematically precise,
Dirac’s notation became universal in physics.
Because Dirac’s formalism is vastly more flexible and computationally powerful.

\section{Axioms of Quantum Mechanics}
\label{aqm}

Today we understand that Dirac’s formalism is rigorously captured by:
\begin{itemize}
   \item rigged Hilbert spaces;
\item spectral theorem for unbounded self-adjoint operators;
\item distribution theory
\end{itemize}
and there is no reason not to use it.

Formally a quantum system is defined by four axioms:
\begin{axiom}[States]
    \label{states}
Physical state is a vector in complex separable Hilbert space $ \mathcal{H} $
\[ |\psi \rangle \in \mathcal{H} .\]
\end{axiom}
\begin{axiom}[Observables]
    \label{observs}
Observables are self-adjoint operators $ \hat A= \hat A^\dagger$ on $\mathcal H$. 
\end{axiom}
\begin{axiom}[Dynamics]
    \label{dynamics}
    There is an observable $ \hat H$ (Hamiltonian),which generates time evolution
    (unitary group):
    \begin{equation}
       \label {evol}
 \hat U(t) = e^{-i \hat Ht/\hbar}.
    \end{equation}
\end{axiom}
    Evolution may be seen as action on observables (Heisenberg picture)
\begin{align}
   \hat A(t)&= \hat U(t) \hat A \hat U^{-1}(t), \label {heis1}\\
    \frac{d \hat A}{d t} &= \frac{i}{\hbar}[\hat H,\hat A].\label{heis2}
\end{align}
    Or as action on states (Schr\"odinger picture)
\begin{align}
    |\psi(t)\rangle&= \hat U^{-1}(t) |\psi \rangle, \label{schr1} \\
    \frac{d |\psi \rangle}{d t} &=-\frac{i}{\hbar} \hat H |\psi \rangle. \label{schr2}
\end{align}
\begin{axiom}[Born rule]
    \label{born-rule}
    \[  \frac{\langle{\psi}|{\phi}\rangle}{\|{\psi}\| \|{\phi}\|}  \]
     is probability amplitude of transition between states.
    Expectation value of an observable in a state 
    \begin{equation}
       \label {exp1}
    \langle \hat A \rangle_\psi = \frac{\langle \psi| \hat A |\psi \rangle} {\langle \psi|\psi \rangle}
    \end{equation}
    comprises all available information.
\end{axiom}

\section{Quantization }
\label{quantization}

Quantization is the collection of procedures by which a classical mechanical 
system is associated with a non-relativistic quantum theory. 
Naturally  quantization is neither unique nor algorithmic in general 
\footnote{ Logically we should speak of de-quantization --- transition from quantum to classical system }, 
and no single scheme applies uniformly to all classical systems.Below  we analyze some quantization schemes, 
discuss no-go theorems limiting the existence of consistent quantization maps.

\subsection {Canonical Quantization}
\label{canonical-quantization}

Canonical quantization assigns self-adjoint operators $\hat q^i, \hat p_j$ on Hilbert space $ \mathcal{H} $
satisfying canonical commutation relations
\begin{equation}
[\hat q^i, \hat p_j] = i\hbar , \delta^i{}_j.
\end{equation}
to classical observables $q^i, p_j$ --- coordinates on phase space $ \mathbb{R}^{2 n} $.
Though this prescription is insufficient for general systems and obscures deeper structural issues: 
the relation between classical observables and quantum operators, the (non-)existence of functorial 
quantization procedures, and the origin of inequivalent quantizations. 

For most practical problems it is realized as
\begin{equation}
q^i \mapsto \hat q^i = q^i, \qquad p_j \mapsto \hat p_j = -i\hbar \frac{\partial}{\partial q^j},
\end{equation}
acting on $\mathcal{H}=L^2(Q,dq)$. For general Hamiltonians, operator ordering ambiguities arise, 
and additional analytic input is required to ensure self-adjointness and well-defined dynamics.
The standard name "canonical quantization" is misleading.

\subsection{Dirac Quantization}
\label{dirac-quantization}

A classical non-relativistic system is described by a symplectic manifold $(M,\omega)$, 
typically $M=T^*Q$ with canonical symplectic form
\begin{equation}
\omega = dq^i \wedge dp_i.
\end{equation}
Observables are smooth functions $f \in C^\infty(M)$ equipped with the Poisson bracket
\begin{equation}
{f,g} = \omega(X_f,X_g),
\end{equation}
where $X_f$ is the Hamiltonian vector field defined by $\iota_{X_f}\omega = df$. 
The resulting Poisson algebra encodes both kinematics and dynamics.

Time evolution is generated by a Hamiltonian function $H$ according to
\begin{equation}
\frac{df}{dt} = {f,H}.
\end{equation}
This structure motivates the canonical correspondence between Poisson brackets and 
commutators in quantum theory.

One often idealizes quantization as a map
\begin{equation}
Q : C^\infty(M) \to \mathcal{A}(\mathcal{H}),
\end{equation}
from classical observables to operators on a Hilbert space $\mathcal{H}$, satisfying 
linearity, reality, normalization, and the Poisson--commutator correspondence
\begin{equation}
Q({f,g}) = \frac{1}{i\hbar}[Q(f),Q(g)].
\end{equation}

However the Groenewold--Van Hove theorem shows that no such map exists on the full 
Poisson algebra of observables for generic phase spaces. 
Even for $M=\mathbb{R}^{2n}$, the correspondence can be consistently imposed only on 
restricted subalgebras \footnote{ In my paper \cite{panfil-q} I studied these subalgebras},
demonstrating that quantization is intrinsically limited.

\subsection{Geometric Quantization}


Geometric quantization begins with a complex line bundle $L \to M$ equipped with a connection $\nabla$ satisfying
\begin{equation}
\mathrm{curv}(\nabla) = \frac{1}{\hbar}\omega.
\end{equation}
This requires the integrality condition $[\omega]/2\pi\hbar \in H^2(M,\mathbb{Z})$.


To obtain a physical Hilbert space, one introduces a polarization selecting a subset of compatible observables. The choice of polarization is generally non-unique and may not exist globally.

\subsection{Deformation Quantization}

Deformation quantization replaces the pointwise product of functions on phase space by a star product
\begin{equation}
f \star g = fg + \frac{i\hbar}{2}{f,g} + O(\hbar^2).
\end{equation}
Quantum mechanics is then viewed as a deformation of the classical observable algebra rather than as an operator theory ab initio.

\subsection{Path Integral Quantization}

In the path-integral approach, transition amplitudes are formally expressed as
\begin{equation}
\langle q_f,t_f | q_i,t_i \rangle = \int \mathcal{D}q, e^{\frac{i}{\hbar}S[q]},
\end{equation}
where $S$ is the classical action. While powerful and intuitive, this formulation typically requires auxiliary constructions for mathematical rigor.

\subsection{Stochastic quantization}
\label{stochastic-quantization}

\subsection{Tomographic Quantization}
\label{tomographic-quantization}

\section{Canonical Quantization Examples}
\label{cqe}
\subsection{Quantum Dynamics of a Wave Packet in a Square Well}

Consider a particle confined to a one-dimensional infinite square well of width $L$,
\begin{equation}
V(x) = 
\begin{cases}
0, & 0 < x < L,\\
\infty, & \text{otherwise}.
\end{cases}
\end{equation}

A \emph{wave packet} is a superposition of stationary states,
\begin{equation}
\Psi(x,0) = \sum_{n=1}^{\infty} c_n \, \phi_n(x), 
\qquad
\phi_n(x) = \sqrt{\frac{2}{L}} \sin \frac{n \pi x}{L},
\end{equation}
where the coefficients $c_n$ encode the initial localization and momentum distribution of the particle.  
Time evolution follows from the Schr\"odinger equation,
\begin{equation}
\Psi(x,t) = \sum_{n=1}^{\infty} c_n \, \phi_n(x) \, e^{-i E_n t/\hbar}, 
\qquad
E_n = \frac{n^2 \pi^2 \hbar^2}{2 m L^2}.
\end{equation}
The quadratic energy spectrum is responsible for the characteristic quantum features described below.

For a wave packet narrowly localized in both position and momentum (centered around a large $n_0$):
\begin{itemize}
    \item The packet moves approximately according to the classical trajectory,
    \begin{equation}
    x_{\mathrm{cl}}(t) = x_0 + v t,
    \end{equation}
    reflecting elastically at the walls.
    \item Ehrenfest's theorem ensures that the expectation values of position and momentum follow classical equations of motion at short times.
\end{itemize}

Due to the quadratic spectrum $E_n \sim n^2$, different stationary components acquire phases at different rates.  
Consequences include:
\begin{itemize}
    \item Spreading (dispersion) of the wave packet.
    \item Destruction of the initial localization by interference.
    \item A characteristic \emph{collapse time} after which the probability density appears irregular.
\end{itemize}
The evolution remains fully unitary; no energy loss or dissipation occurs.

The exact quadratic spectrum leads to wave packet \emph{revivals}:
\begin{equation}
T_{\mathrm{rev}} = \frac{4 m L^2}{\pi \hbar}.
\end{equation}
At $t = T_{\mathrm{rev}}$, the wave packet reconstructs itself up to a global phase,
\begin{equation}
\Psi(x,T_{\mathrm{rev}}) = e^{i \theta} \Psi(x,0),
\end{equation}
while at rational fractions $t = \frac{p}{q} T_{\mathrm{rev}}$, \emph{fractional revivals} occur, splitting the packet into $q$ localized copies.

Inside the infinite well:
\begin{itemize}
    \item A self-adjoint momentum operator compatible with the boundary conditions does not exist.
    \item Operationally, momentum must be described via a POVM, not a PVM.
    \item The packet behaves as if momentum reverses sign at the walls, producing strong interference near the boundaries.
\end{itemize}

The dynamics of a wave packet in a square well illustrate:
\begin{description}
    \item [Short-time classical motion:] The packet moves and reflects approximately like a classical particle.
    \item [Quantum dispersion:] Interference between stationary components spreads the packet.
    \item [Exact revivals:] The wave packet periodically reconstructs itself due to the quadratic spectrum.
    \item [Unitary boundary-induced interference:] Reflections at the walls generate characteristic oscillations without dissipation.
\end{description}
This system exemplifies the interplay between classical intuition and fundamentally quantum effects.

\section{The Ground State Determines  Dynamics }
\label{gsd}

In classical mechanics, dynamics is specified by a Hamiltonian function; individual solutions do not determine the Hamiltonian. In quantum mechanics, however, the structure of the Schr\"odinger equation implies a striking inversion: knowledge of a single, sufficiently regular stationary state---specifically the nondegenerate ground state---fixes the potential and thereby all dynamics.

This observation appears in scattered forms across the literature: inverse Sturm--Liouville problems, the Riccati formulation of one-dimensional quantum mechanics, supersymmetric factorizations, and the Hohenberg--Kohn theorem. Yet its conceptual significance is rarely stated explicitly. The aim of this paper is to present a clean, general argument and to clarify its physical and philosophical implications.

Consider a single nonrelativistic particle of mass $m$ on $\mathbb{R}^d$ governed by
\begin{equation}
\hat H = -\frac{\hbar^2}{2m}\nabla^2 + V(\mathbf{x}),
\end{equation}
with $V$ real-valued, locally bounded, and confining enough to admit a normalizable ground state $\psi_0>0$ almost everywhere.

We assume:
\begin{enumerate}
\item The ground state is nondegenerate.
\item $\psi_0(\mathbf{x})$ is twice differentiable and nowhere vanishing.
\item Boundary conditions (e.g. decay at infinity) are known.
\end{enumerate}
These assumptions are standard in both mathematical physics and quantum chemistry.

Let $E_0$ be the ground-state energy. The time-independent Schr"odinger equation reads
\begin{equation}
-\frac{\hbar^2}{2m}\nabla^2\psi_0 + V\psi_0 = E_0\psi_0.
\end{equation}
Solving for the potential yields
\begin{equation}
    \label{pot1d}
V(\mathbf{x}) = E_0 + \frac{\hbar^2}{2m}\frac{\nabla^2\psi_0(\mathbf{x})}{\psi_0(\mathbf{x})}.
\end{equation}
Thus, given $\psi_0$ and $E_0$, the potential is uniquely determined pointwise. Since adding a constant to $V$ merely shifts all energies uniformly, the potential is fixed up to an additive constant even if $E_0$ is unknown.

The positivity of $\psi_0$ ensures that the reconstruction formula is well defined. Any other potential $V'$ admitting the same ground state must coincide with $V$ almost everywhere. Hence the ground state uniquely specifies the Hamiltonian.

Once $V$ is fixed, the entire spectral problem is fixed: all excited eigenstates and eigenvalues are determined, the propagator $U(t)=e^{-i\hat H t/\hbar}$ is fixed, and all expectation values and transition amplitudes follow. Therefore, the ground state determines not merely low-energy physics but the complete time evolution of the system.

In one dimension, the result is a special case of inverse Sturm--Liouville theory, where knowledge of the lowest eigenfunction suffices to reconstruct the potential locally.

Defining the superpotential
\begin{equation}
W(x) = -\frac{\hbar}{\sqrt{2m}}\frac{\psi_0'(x)}{\psi_0(x)},
\end{equation}
we obtain
\begin{equation}
V(x) = W^2(x) - \frac{\hbar}{\sqrt{2m}}W'(x) + E_0,
\end{equation}
showing again that $\psi_0$ encodes the full Hamiltonian structure.

The Hohenberg--Kohn theorem states that the ground-state density uniquely determines the external potential up to a constant. The present result may be viewed as a wavefunction-level strengthening: the full ground-state amplitude determines the Hamiltonian directly, without recourse to variational principles.

Let the ground-state wavefunction be written as
\begin{equation}
\psi_0(\mathbf{x}) = e^{-\phi(\mathbf{x})},
\end{equation}
where $\phi$ is a smooth real-valued function. Positivity of the ground state guarantees that this representation is globally well defined. Substituting into the reconstruction formula yields
\begin{equation}
V(\mathbf{x}) = E_0 + \frac{\hbar^2}{2m}\Big[(\nabla \phi)^2 - \nabla^2 \phi\Big].
\end{equation}
The potential is thus a scalar functional constructed from first and second derivatives of $\phi$, suggesting an intrinsic geometric interpretation.

Define a conformally flat metric on configuration space
\begin{equation}
g_{ij}(\mathbf{x}) = e^{-2\phi(\mathbf{x})},\delta_{ij}.
\end{equation}
The associated Laplace--Beltrami operator satisfies
\begin{equation}
\Delta_g = e^{2\phi}\bigl(\nabla^2 - (d-2)\nabla \phi \cdot \nabla\bigr).
\end{equation}
In this formulation, the ground-state Schr\"odinger equation may be reinterpreted as a geometric condition relating curvature invariants of $g_{ij}$ to the energy $E_0$. The classical potential thus emerges as a curvature scalar of a geometry induced by the quantum state.

Introduce the quantum potential
\begin{equation}
Q[\psi] = -\frac{\hbar^2}{2m}\frac{\nabla^2\psi}{\psi}.
\end{equation}
For the ground state one has
\begin{equation}
V = E_0 - Q[\psi_0].
\end{equation}
Thus the classical potential coincides with minus the quantum potential evaluated on the ground state. From a Hamilton--Jacobi viewpoint, $\phi$ plays the role of an imaginary action, and the ground state defines a preferred foliation of configuration space.

Physical quantum states correspond to rays in Hilbert space, forming a projective manifold equipped with the Fubini--Study metric. The ground state selects a distinguished point in this space. The reconstruction of $V$ from $\psi_0$ may be interpreted as lifting geometric data from projective Hilbert space to a differential-geometric structure on configuration space, reinforcing the view that the Hamiltonian is not primitive but induced by state geometry.

Degenerate ground states may lead to non-uniqueness. Excited states cannot be used directly due to the presence of nodes. For interacting many-body systems the statement generalizes formally, but reconstruction becomes nonlocal in configuration space. The argument does not extend straightforwardly to relativistic or quantum field theoretic settings.

This determinacy challenges the intuition that dynamics is prior to states. In the Schr\"odinger framework, the ground state contains the full dynamical content. The Hamiltonian may therefore be regarded as a derived object, secondary to a privileged stationary state. This perspective aligns naturally with structural and geometric interpretations of quantum theory.


\subsection{Explicit Example}

Consider the nontrivial ground state
\begin{equation}
    \label{gsx4}
\psi_0(x) = A \exp\Big(-\frac{a}{4} x^4\Big), \quad a>0.
\end{equation}

Substituting \eqref{gsx4} into \eqref{pot1d}, we compute:
\begin{align*}
\psi_0'(x) &= -a x^3 \psi_0(x), \
\psi_0''(x) &= (a^2 x^6 - 3 a x^2) \psi_0(x).
\end{align*}
Thus the potential is
\begin{equation}
V(x) = E_0 + \frac{\hbar^2}{2m} (a^2 x^6 - 3 a x^2),
\end{equation}
a sextic potential with a negative quadratic term, uniquely determined by $\psi_0$.

Writing excited states as $\psi_n = \psi_0 P_n$, the Sturm--Liouville operator for $P_n$ is
\begin{equation}
L = -\frac{\hbar^2}{2m} \left( \frac{d^2}{dx^2} + 2 a x^3 \frac{d}{dx} \right),
\end{equation}
with weight $w(x) = \psi_0^2 = A^2 e^{- \frac{a}{2} x^4}$. The eigenvalue equation is
\begin{equation}
L P_n = (E_n - E_0) P_n,
\end{equation}
which defines the spectrum and orthogonal polynomials (Freud polynomials) relative to $w(x)$.
Numerical or recursive methods can be used to compute $E_n$ and $P_n$.

